% Section 4.3, from lectures/L04/appendix/b_exercises.md.

\section{Exercises}
\label{c4:sec:exercises}\label{c4:app:b}

\begin{aside}
Nine, ending with the Cross-check. Do them in order; Exercise~\ref{c4:ex:6} onwards need the modules
from the lab.

Where an exercise can be checked mechanically, a \textbf{Check yourself} line says how. Several have
a plausible wrong answer that is worth walking into before you find it.

Everything here runs against the target from \code{make boot}, or against the built tree under
\file{build/kernel}. Nothing needs the network.

\textbf{Solutions.} The exercises that are not Code are answered in \secref{sol:sec:c4}. The Code
exercises have no solution: each is checked by running it, as its Check yourself line says.
\end{aside}

% ------------------------------------------------------------------------------------------
\exercise[fillaccent2]{A module is not a program}{Recall}
\label{c4:ex:1}

\xpart{a} \code{qa_hello_init} returns 0. What is running afterwards, and what is resident?

\xpart{b} Why is it a mistake to put a \code{while (1)} loop in a module's init function? What
should you have written instead, and what happens to \code{insmod} if you do it anyway?

\xpart{c} What do \code{__init} and \code{__exit} actually do? Neither is a hint to the reader.

\xpart{d} The boot log says \code{Freeing unused kernel memory: 4672K}. What relationship does that
message have to \code{__init}?

\xpart{e} Name the four \code{MODULE_*} macros from the lab's module, and say which one is not
optional and what happens if you leave it out.

% ------------------------------------------------------------------------------------------
\exercise[fillaccent2]{Four commands}{Recall}
\label{c4:ex:2}

\xpart{a} State in one sentence each what \code{insmod}, \code{modprobe}, \code{depmod} and
\code{rmmod} do.

\xpart{b} \code{insmod} and \code{modprobe} are often described as the same command. Give the
concrete difference, and describe the experiment that demonstrates it using the lab's
\code{qa_export} and \code{qa_user}.

\xpart{c} What file does \code{modprobe} read to decide what to load first, and what generates it?

\xpart{d} \code{rmmod qa_export} fails while \code{qa_user} is loaded. Where, in \file{/sys}, can
you see \emph{which} module is holding it, rather than merely how many are?

% ------------------------------------------------------------------------------------------
\exercise[fillaccent3]{Reading \code{vermagic}}{Hand calculation}
\label{c4:ex:3}

\code{modinfo} on the lab's module reports:

\begin{console}
vermagic: 6.12.30 SMP preempt mod_unload modversions aarch64
\end{console}

\xpart{a} Name what each of the six fields is asserting about the kernel this module may be loaded
into.

\xpart{b} For each of the following changes to the kernel, say whether this module would still load,
and why:
\begin{itemize}
  \item The kernel is rebuilt from the same source with no config change.
  \item The kernel is rebuilt with \code{\# CONFIG_SMP is not set} in \file{kernel/local.config}.
    Read \file{arch/arm64/Kconfig} before answering.
  \item The kernel is upgraded from 6.12.30 to 6.12.31.
  \item The kernel is rebuilt with \code{CONFIG_PREEMPT_RT=y} instead of \code{CONFIG_PREEMPT=y}.
\end{itemize}

\xpart{c} \code{modversions} appears in the list. What does \code{CONFIG_MODVERSIONS} add beyond the
version string, and what problem does it solve that \code{vermagic} alone does not?

\xpart{d} Explain why a mismatch is reported as an error rather than being allowed with a warning.
What is the failure mode that policy is protecting you from?

% ------------------------------------------------------------------------------------------
\exercise[fillaccent3]{The licence, at two stages}{Hand calculation}
\label{c4:ex:4}

A module declares \code{MODULE_LICENSE("Proprietary")} and calls a function exported with
\code{EXPORT_SYMBOL_GPL}.

\xpart{a} You build it in the same directory as the module that exports the symbol. What happens, at
which build step, and what exactly does the message say?

\xpart{b} You build it in a directory of its own, where the build cannot see the exporting module's
symbols. What happens now, and why is it a \emph{different} message?

\xpart{c} Neither of those produced a \code{.ko}. Describe the circumstances in which a proprietary
module using a GPL-only symbol does build successfully and then fails at \code{insmod}.

\xpart{d} In that case, what does \code{insmod} print? Read \secref{c4:app:a5} and say why the
message does not mention a licence, in terms of what the module loader actually does with GPL-only
symbols.

\xpart{e} Why is that the worst of the three messages to receive?

% ------------------------------------------------------------------------------------------
\exercise[fillaccent4]{Choosing a log level}{Design}
\label{c4:ex:5}

For each, choose from \code{pr_debug}, \code{pr_info}, \code{pr_warn}, \code{pr_err}, and justify it
in one sentence.

\xpart{a} The driver has probed successfully and is ready.

\xpart{b} A device returned a value outside its documented range, and you clamped it.

\xpart{c} \code{ioremap} returned NULL and \code{probe} is about to fail.

\xpart{d} You want to print the value of every register during development.

\xpart{e} The device tree specified a clock rate the hardware cannot produce, so you used the
nearest.

Then: \textbf{f)} why is \code{pr_debug} not simply \code{pr_info} that you comment out before
shipping? Give two reasons, one about the binary and one about the field.

% ------------------------------------------------------------------------------------------
\exercise{Your first module}{Code}
\label{c4:ex:6}

Write \file{lectures/L04/lab/qa_hello.c}. It must:
\begin{itemize}
  \item Print a greeting from its init function and a farewell from its exit function.
  \item Take a string parameter \code{who}, defaulting to \code{world}, readable but \textbf{not}
    writable via sysfs.
  \item Take an integer parameter \code{count}, defaulting to 1, readable \textbf{and} writable via
    sysfs, that controls how many times the greeting is printed.
  \item Carry \code{MODULE_LICENSE}, \code{MODULE_AUTHOR}, \code{MODULE_DESCRIPTION} and
    \code{MODULE_VERSION}.
  \item Use \code{pr_fmt} so every message is prefixed with the module name, and contain one
    \code{pr_debug} that says something only useful when switched on.
\end{itemize}

\checkyourself \code{insmod ./qa_hello.ko who=embedded count=2} prints the greeting twice; \code{cat
/sys/module/qa_hello/parameters/who} returns \code{embedded}; writing to that file fails with
\code{Permission denied}, and writing to \code{count} succeeds.

% ------------------------------------------------------------------------------------------
\exercise{One module using another}{Code}
\label{c4:ex:7}

Write \file{qa_export.c}, which exports \code{qa_answer()} returning 42 with
\code{EXPORT_SYMBOL_GPL}, and \file{qa_user.c}, which calls it from its init function and prints the
result.

\xpart{a} Load them in the wrong order deliberately. Record the exact error.

\xpart{b} Load them in the right order, then try to \code{rmmod} them in the wrong order. Record
what happens, and find the holder in \file{/sys}.

\xpart{c} \file{qa_export.c} also exports \code{qa_answer_open()} with plain \code{EXPORT_SYMBOL}.
Write a third module declaring \code{MODULE_LICENSE("Proprietary")} that calls \textbf{only} that
one. Does it build? Does it load?

\xpart{d} Read \file{/proc/sys/kernel/tainted} before loading anything, after loading
\code{qa_export}, and after loading the proprietary one. You get three different numbers. Decompose
each into bits, and name the two taints involved using \file{include/linux/panic.h}. Then look at
the flags in \file{/proc/modules} and in \code{/sys/module/<name>/taint}: which letters appear
against which module, and why does the proprietary module carry two?

\xpart{e} Loading the proprietary module prints one further line that is not about licensing at all:

\begin{console}
Disabling lock debugging due to kernel taint
\end{console}

\noindent Find out what was disabled. Then say what that means for Chapter~\ref{c7:ch}, which is
built entirely on that facility, and state the rule it implies for a development kernel.

\checkyourself \code{make test L=L04} passes, with twelve checks. Note that its check on
\code{rmmod} deliberately asserts on \code{lsmod} rather than on \code{rmmod}'s exit status;
Exercise~\ref{c4:ex:8} is why.

% ------------------------------------------------------------------------------------------
\exercise[fillaccent4]{A test that cannot fail}{Design}
\label{c4:ex:8}

The lab's \file{run.sh} checks that a module in use cannot be unloaded. It does so by attempting the
\code{rmmod} and then looking at \code{lsmod}, rather than by checking \code{rmmod}'s exit status.

\xpart{a} Run \code{rmmod qa_export} while \code{qa_user} is loaded, then immediately \code{echo
$?}. What do you get?

\xpart{b} Given that, what would a test written the obvious way conclude, on a kernel that correctly
refused \emph{and} on one that wrongly allowed it?

\xpart{c} Name the general rule this is an instance of, and give one other example from
\secref{c2:app:a8} of the same class of problem.

% ------------------------------------------------------------------------------------------
\exercise[fillaccent]{How big is a module}{Cross-check}
\label{c4:ex:9}

Three measurements of the same module, differing by a factor of five. Predict, measure, reconcile.

\xpart{a} \textbf{The file.} With the lab built, record the size of \file{qa_hello.ko} on disk.

\xpart{b} \textbf{The content.} Run \code{size -A} on it, in the container so you get the cross
toolchain's version:

\begin{shell}
aarch64-linux-gnu-size -A lectures/L04/lab/qa_hello.ko
\end{shell}

\noindent Record the \code{Total}. Then record the sum of only those sections that will actually be
\textbf{loaded}: leave out \code{.modinfo}, \code{__versions} and \code{.comment}, which are
metadata the loader reads and does not keep, and \code{.symtab} and \code{.strtab}, of which it
keeps only a small copy for \file{/proc/kallsyms}. The \code{.note.*} sections are kept, and are
small.

\xpart{c} \textbf{Predict.} From \textbf{b)}, write down what you expect the module to occupy in
kernel memory once loaded.

\xpart{d} \textbf{Measure.} Boot, load the module, and read its actual size two ways:

\begin{shell}
grep qa_hello /proc/modules
cat /sys/module/qa_hello/coresize
\end{shell}

\xpart{e} \textbf{Reconcile.} Answer each:
\begin{itemize}
  \item By what factor does the measured size exceed your prediction?
  \item The measured number is round. What is it a multiple of, and how many of them?
  \item Load \code{qa_export} and compare \emph{its} \code{coresize} with \code{qa_hello}'s. They
    are not the same multiple. Before reading further, say what would have to differ between two
    modules for one to need more of these units than the other.
  \item Now find out why. List the section addresses and sort them:
\begin{shell}
for f in /sys/module/qa_hello/sections/.*; do
    echo "$(cat "$f" 2>/dev/null) $(basename "$f")"
done | sort
\end{shell}

    The addresses fall into distinct groups. How many groups, what is in each, and what is the
    spacing between them? What does that tell you about how a module is laid out in memory?
  \item \file{include/linux/module.h} defines \code{enum mod_mem_type}. How many of its entries are
    \emph{core} types rather than init types, and how does that number relate to what you measured?
    The allocation itself is one line in \file{kernel/module/main.c}; find it.
  \item Why is the module laid out this way rather than packed into one contiguous region? The
    answer is about permissions, and it is the same reason the kernel's own \code{.text} and
    \code{.data} are separate.
\end{itemize}

\xpart{f} \file{/sys/module/qa_hello/initsize} reads 0 after the module has finished loading, even
though the module plainly has an init function. What happened to it, and what does that have to do
with Exercise~\ref{c4:ex:1}(d)?

\xpart{g} Someone asks how much memory your driver uses. Give a single number and a sentence, and
say which of your three measurements you chose and why. Then say what your answer would be if the
driver also allocated a 64 KB buffer with \code{kmalloc} in its init function, and which of the
three numbers would have changed.
